\subsection{NN-G027: nonconstant ratios produce smaller explicit realizations}
\paragraph{Statement.}
There is a closed-table effective associative dimension ring
with $3452$ cancelling pairs, of which $2852$ are in the
bottom cell. The addendum below improves these numbers to
$3352$ and $2752$, respectively. Both constructions have
nonconstant rational functions
$G_i/(W\delta_i)$, so it lies outside the exact $4520$
class of NN-G026. Thus the constant-ratio restriction is
substantive. No exact optimum over the enlarged rational-function
class is asserted.

\paragraph{Construction and associativity.}
Put
\[
 \sigma=D(D-3),\qquad h=\sigma+1,\qquad B=5\sigma+4.
\]
Choose
\[
 \begin{aligned}
 \delta_1&=\sigma+2,&\delta_2&=3\sigma+2,&W&=3D,\\
 G_1&=6Dh,&G_2&=DB,&S&=2DhB,\\
 U_1&=3D,&U_2&=D,&V_1&=15D-6,&V_2&=2D^2+D,\\
 A_1&=6D^2h-B(5D-2),&
 A_2&=D^2B-2h(2D^2+D),&
 \Lambda_i&=A_i/\delta_i.
 \end{aligned}                                                       \tag{G027.1}
\]
The actual within-cell coefficients involving $\Lambda_i$
are $A_i$ multiplied by nonnegative powers of $D$.
They are integral polynomials even though the auxiliary
$\Lambda_i$ are rational functions.

All three scalar-block identities follow directly:
\[
 \frac{G_1G_2}{S}=3D=W,\qquad
 G_1-D\delta_1=G_2,\qquad
 G_2-D\delta_2=G_1/3.
\]
Consequently $G_i(G_i-D\delta_i)/S=U_i$.
The two definitions of $A_i$ give
\[
 \frac{G_1(DG_1-A_1)}S=3(5D-2)=V_1,\qquad
 \frac{G_2(DG_2-A_2)}S=V_2.
\]
Also $\Lambda_i(D^2-\Lambda_i)\ne0$: the leading
terms of $\Lambda_1,\Lambda_2$ are respectively
$6D^2$ and $D^2/3$. The injective generic block
construction of NN-G011 therefore proves associativity
over $\mathbb Q(D)$. As all actual structure constants
are integral, its identities hold over $\mathbb Z[D]$.
This argument makes no semisimplicity claim at the
special values where an auxiliary factor has a pole.

\paragraph{Coefficient verification and exact score.}
The following symmetric Laurent lists give every base
coefficient needed for the effectivity, domination and
parity checks:
\[
\begin{array}{c|l|l}
 &\text{dimension coefficient}&\text{absolute trace target}\\\hline
 \delta_1&(2,1,1)&(0,1,1)\\
 \delta_2&(2,3,3)&(0,1,1)\\
 A_1&(14,5,1,5,6)&(2,1,1,1)\\
 A_2&(2,3,4,3,1)&(2,3,2,1,1)\\
 G_1&(24,24,18,6)&(16,12,4)\\
 G_2&(18,19,15,5)&(14,11,5,1)\\
 S&(172,152,114,78,40,10)&(144,120,68,24,4)\\
 V_1&(24,15)&(6,5)\\
 V_2&(14,9,2)&(6,3,2)
\end{array}                                                         \tag{G027.2}
\]
The remaining base coefficients are $W=U_1=3D$ and
$U_2=D$. Every entry in the middle column dominates
the right column coefficientwise with the same parity.
The closed multiplication table only multiplies these
base coefficients by nonnegative powers of $D$;
for signed traces use the coefficientwise triangle
inequality. Hence all $132$ positions satisfy the
required conditions. The included independent
integer-Laurent verifier also checks all $1000$
associativity triples and reports zero failures.

At $D=4$ the bottom mass is
\[
 960+20(120+96)+200\cdot12+20(12+4+54+36)=9800.
\]
The two cell characters have masses $6,14$, and the
actual $A_1,A_2$ have masses $48,24$. Thus their
contribution is $26(6+14)+14(48+24)=1528$.
The total is $T=11328$, yielding
\[
 (11328-4424)/2=3452,\qquad (9800-4096)/2=2852.
 \tag{G027.3}
\]
Finally
\[
 \frac{G_1}{W\delta_1}=\frac{2(\sigma+1)}{\sigma+2},
 \qquad
 \frac{G_2}{W\delta_2}=\frac{5\sigma+4}{3(3\sigma+2)},
\]
which are both nonconstant. The lower score is therefore
consistent with NN-G026, and proves that extending
its optimum beyond constant ratios would be false.
\hfill$\square$

\paragraph{A verified refinement to 3352 pairs.}
The same identities permit a smaller realization. Keep
$H=D^2-3D+1$, and now put
\[
 J=5D^2-17D+10.
\]
Define
\[
 \begin{aligned}
 \delta_1&=6H-J=D^2-D-4,&
 \delta_2&=J-2H=3D^2-11D+8,\\
 G_1&=6DH,&G_2&=DJ,&S&=2DHJ,\\
 W&=3D,&U_1&=3D,&U_2&=D,\\
 V_1&=15D,&V_2&=2D^2+D-4,\\
 A_1&=6D^2H-5DJ,&
 A_2&=D^2J-2H(2D^2+D-4),&\Lambda_i&=A_i/\delta_i.
 \end{aligned}                                                       \tag{G027.4}
\]
Again $G_1-D\delta_1=G_2$ and
$G_2-D\delta_2=G_1/3$, while the definitions of $A_i$
give the two $V_i$ identities. Neither
$\Lambda_i(D^2-\Lambda_i)$ is identically zero, so the
same injective block construction proves associativity.
The required base coefficients are
\[
\begin{array}{c|l|l}
 &\text{dimension coefficient}&\text{absolute trace target}\\\hline
 \delta_1&(0,3,1)&(0,1,1)\\
 \delta_2&(4,1,3)&(0,1,1)\\
 A_1&(6,5,1,5,6)&(2,1,1,1)\\
 A_2&(6,5,6,1,1)&(2,3,2,1,1)\\
 G_1&(24,24,18,6)&(16,12,4)\\
 G_2&(18,17,13,5)&(14,11,5,1)\\
 S&(156,140,106,70,36,10)&(144,120,68,24,4)\\
 V_1&(30,15)&(6,5)\\
 V_2&(10,9,2)&(6,3,2)
\end{array}                                                         \tag{G027.5}
\]
They prove every base condition; powers of $D$ propagate
them to the full table. The independent verifier
\texttt{verify\_dimension\_3352.py} checks all $132$
positions and all $1000$ triples with zero failures.

The bottom mass is
\[
 880+20(120+88)+200\cdot12+20(12+4+60+32)=9600.
\]
The cell mass is $26(8+12)+14(40+32)=1528$.
Consequently
\[
 T=11128,\qquad
 (11128-4424)/2=3352,\qquad(9600-4096)/2=2752.
 \tag{G027.6}
\]
The two ratios are $2H/(6H-J)$ and
$J/(3(J-2H))$, both nonconstant. This strengthens the
explicit upper bound to $3352$, retaining the original
$3452$ example as the first checked construction.
NN-G028 proves an exact optimum for the indicated
fixed-bottom subclass; no unrestricted optimum is
claimed here.\hfill$\square$

\paragraph{Geometric meaning and residual scope.}
The actual polynomials in (G027.2) are effective
Artin--Tate multiplicity objects and define another
NN-G010 finite-kernel monoidal realization of the
same trace table. With its chosen generators, the
sum over multiplication positions and grading pieces
of half the excess of dimension over absolute signed
trace is $3452$ for the first construction, and $3352$
for (G027.4), with bottom-cell contributions $2852$
and $2752$, respectively.
This is an explicit upper bound for a presentation-dependent
quantity, not a universal invariant of all categories
realizing the abstract algebra.

The rational $\Lambda_i$ and the two nonconstant ratios
are generic coordinate factorizations; no object of
fractional dimension or meromorphic sheaf is inferred.
Only the verified actual coefficients define objects.
The precise dictionary is SGA 3, I, 4.7.2--4.7.3 and
archive cote 2, scan 17 through NN-G010. The example
leaves the exact optimum for arbitrary rational ratios
and for other table shapes open. It does not give a
canonical family, a rank-four basis, PH0/PH1 or
nonstandard purity.

